\documentclass[12pt]{article}
\usepackage{amsmath,amssymb,amsfonts,amsthm}
\newtheorem*{mainresult}{Main Result}
\usepackage{hyperref}
\usepackage{geometry}
\usepackage{authblk}
\geometry{margin=1in}

\newtheorem{theorem}{Theorem}[section]
\newtheorem{lemma}[theorem]{Lemma}
\newtheorem{proposition}[theorem]{Proposition}
\newtheorem{corollary}[theorem]{Corollary}
\theoremstyle{definition}
\newtheorem{definition}[theorem]{Definition}
\newtheorem{assumption}[theorem]{Assumption}
\theoremstyle{remark}
\newtheorem{remark}[theorem]{Remark}

\newcommand{\Hthree}{\mathbb{H}^3}
\newcommand{\R}{\mathbb{R}}
\newcommand{\Z}{\mathbb{Z}}
\newcommand{\bdry}{\partial\mathbb{H}^3}
\newcommand{\Q}{\mathbb{Q}}
\newcommand{\C}{\mathbb{C}}
\newcommand{\vp}{\varphi}

\title{Weight Spectra of Arithmetic Hyperbolic Lattices\\
and the Mahler Measure of the Alexander Polynomial}
\date{}
\author[1]{Marvin L.\ Gentry}
\affil[1]{Independent Researcher\\
  \texttt{drmlgentry@protonmail.com}\\
  ORCID: 0009-0006-4550-2663}

\begin{document}
\maketitle

\begin{abstract}
Let $\Gamma$ be a torsion-free cocompact lattice in $\mathrm{PSL}(2,\C)$
acting on $\Hthree$ by orientation-preserving isometries, and let
$q(\gamma)=\kappa\,d(0,\gamma\cdot 0)$ denote the logarithmic weight of
$\gamma\in\Gamma$.  We prove three structural results holding for any such
$\Gamma$: (1)~each boundary direction admits a unique minimal-weight
representative; (2)~the weight spectrum has a uniform positive gap
$\Delta=\kappa\,\mathrm{sys}(\Gamma)$ equal to $\kappa$ times the systole;
(3)~the integer coordinates arising from the Milnor--\v{S}varc
quasi-isometry are stable under deformations smaller than half the minimum
separation of boundary directions.

For the arithmetic lattice $\Gamma=\pi_1(M_{\mathrm{PMNS}})$ arising from
the hyperbolic 3-manifold $M_{\mathrm{PMNS}}=\mathrm{OrientableClosedCensus}[1]$
(the second-smallest closed hyperbolic 3-manifold, volume $0.9814$,
$H_1=\Z/5$), we compute the weight spectrum explicitly and identify
its gap with the mass-lattice spacing of the Standard Model fermion
masses.  The key arithmetic invariant is the Alexander polynomial of
the cusped ancestor \texttt{m003}, which equals $\Delta(t)=t^2+3t+1$
and has roots $-\vp^2$ and $-\vp^{-2}$, where $\vp=(1+\sqrt{5})/2$
is the golden ratio.

\begin{mainresult}[Algebraic Bridge, informal]
The Mahler measure of the Alexander polynomial of the cusped ancestor of
$M_\mathrm{PMNS}$ satisfies
\[
\log M(\Delta_{M_\mathrm{PMNS}}) = 2\log\vp = 2\,\mathrm{Reg}(\Q(\sqrt{5})),
\]
which equals twice the uniform gap in the weight spectrum of the Standard
Model mass lattice $\{m_e\cdot\vp^{q/4}:q\in\Z\}$.  This provides a
purely topological origin for the golden ratio in the lepton sector,
independent of any numerical fit.
\end{mainresult}

All manifold invariants (Alexander polynomial, adjoint torsion, Chern--Simons
invariant, length spectrum) are computed via \texttt{Sage~10+SnapPy~3.3}.
\end{abstract}

\section{Introduction}

Discrete logarithmic hierarchies appear in warped extra dimensions,
renormalization group flows, and bulk-boundary correspondences in
holography.  A physically useful framework must explain three properties:
(i)~why each weight level supports a unique dominant contribution,
(ii)~why subleading contributions are suppressed by a definite amount,
and (iii)~why the organization is stable under perturbations.
In this paper we prove that all three properties are consequences of the
geometry of a cocompact lattice action on hyperbolic space, and we
identify the suppression gap in a concrete arithmetic example with the
regulator of a quadratic number field.

The motivating example is the \emph{PMNS manifold}
$M_\mathrm{PMNS}$~\cite{Gentry2026masses}, the closed hyperbolic
3-manifold of second-smallest volume ($0.9814$) with first homology
$H_1=\Z/5$.  Its fundamental group $\pi_1(M_\mathrm{PMNS})$ is a
torsion-free cocompact lattice in $\mathrm{PSL}(2,\C)$.
The cusped ancestor \texttt{m003} (the $(-2,3,7)$ pretzel knot
complement) has Alexander polynomial $\Delta(t)=t^2+3t+1$, whose
roots are $-\vp^2$ and $-\vp^{-2}$ (Theorem~\ref{thm:alexander}).
The Mahler measure $M(\Delta)=\vp^2$ determines the systole of
the associated weight lattice through the algebraic bridge
(Theorem~\ref{thm:bridge}).

\paragraph{Relation to the figure-eight knot.}
The figure-eight knot complement (\texttt{m004}) shares the same
hyperbolic volume $2.02988\ldots$ as \texttt{m003} but has Alexander
polynomial $\Delta_{4_1}(t)=t-3+t^{-1}$ with Mahler measure
$3+2\sqrt{2}\approx5.828$, which carries no connection to the golden
ratio.  The two manifolds are not isometric (verified
computationally~\cite{SnapPy}); the volume coincidence does not imply
arithmetic equivalence.

\paragraph{Organisation.}
Section~\ref{sec:setting} sets up the geometric framework.
Sections~\ref{sec:uniqueness}--\ref{sec:rigidity} prove the three
structural results.
Section~\ref{sec:arithmetic} specialises to the arithmetic lattice
$\pi_1(M_\mathrm{PMNS})$ and proves the Algebraic Bridge.
Section~\ref{sec:gap} identifies the suppression gap with the
mass-lattice spacing.

\section{Geometric Setting}
\label{sec:setting}

\begin{assumption}\label{ass:action}
$\Gamma$ is torsion-free, acting properly discontinuously and cocompactly
on $\Hthree$ by orientation-preserving isometries.
Fix basepoint $0\in\Hthree$.
\end{assumption}

\noindent For $\gamma\in\Gamma$ set $q(\gamma)=\kappa\,d(0,\gamma\cdot0)$,
$\kappa>0$.  The \emph{weight spectrum} of $\Gamma$ is
$\mathcal{W}(\Gamma,\kappa)=\{q(\gamma):\gamma\in\Gamma\}$.

\begin{definition}[Isometry types]\label{def:types}
An isometry $\delta\in\mathrm{Isom}(\Hthree)$ is \emph{elliptic} if it
fixes a point in $\Hthree$; \emph{parabolic} if it fixes exactly one
boundary point; \emph{hyperbolic} if it fixes exactly two boundary points.
\end{definition}

\begin{lemma}[No parabolics in cocompact lattices]\label{lem:parabolic}
A cocompact lattice $\Gamma\subset\mathrm{Isom}(\Hthree)$ contains no
parabolic elements~\cite{ratcliffe}.
\end{lemma}

\section{Uniqueness of Minimal Representatives}
\label{sec:uniqueness}

\begin{definition}\label{def:minimal}
$\gamma\in\Gamma$ is a \emph{minimal representative} for
$\xi\in\bdry$ if the geodesic ray from $0$ through $\gamma\cdot0$
has ideal endpoint $\xi$ and $q(\gamma)$ is minimal among all such
elements.
\end{definition}

\begin{theorem}[Uniqueness]\label{thm:uniqueness}
Every $\xi\in\bdry$ admits at most one minimal representative in $\Gamma$.
\end{theorem}

\begin{proof}
Let $\gamma_1,\gamma_2$ be two minimal representatives for $\xi$.
Then $\delta:=\gamma_1^{-1}\gamma_2$ fixes $\xi\in\bdry$, so $\delta$
is elliptic or parabolic.  Since $\Gamma$ is cocompact it contains no
parabolics (Lemma~\ref{lem:parabolic}).  Hence $\delta$ is elliptic,
and since $\Gamma$ is torsion-free, $\delta=\mathrm{id}$.
\end{proof}

\section{The Uniform Suppression Gap}
\label{sec:gap_section}

\begin{definition}
The \emph{systole} of $\Gamma$ is
$\mathrm{sys}(\Gamma)=\min_{\gamma\neq\mathrm{id}}d(0,\gamma\cdot0)>0$.
\end{definition}

\begin{lemma}[Positive minimum weight]\label{lem:gap}
The weight spectrum $\mathcal{W}(\Gamma,\kappa)$ has a positive minimum
value $\Delta:=\kappa\,\mathrm{sys}(\Gamma)>0$, where
$\mathrm{sys}(\Gamma)=\min_{\gamma\neq\mathrm{id}}d(0,\gamma\cdot0)$.
For any $R>0$, the set $\{\gamma\in\Gamma:q(\gamma)\le R\}$ is finite.
\end{lemma}

\begin{proof}
The orbit $\Gamma\cdot0$ is discrete by proper discontinuity and
$\Gamma$ has no fixed points in $\Hthree$ (torsion-free), so
$\mathrm{sys}(\Gamma)>0$.  Every group element satisfies
$q(\gamma)\ge\Delta$.  Finiteness follows because
$\{q(\gamma)\le R\}=\Gamma\cap\overline{B}(0,R/\kappa)$, which is
finite by proper discontinuity.
\end{proof}

\begin{remark}
Two distinct elements may have the same weight (e.g.\ $\gamma$ and
$\gamma^{-1}$ have the same displacement from $0$).
The lemma asserts only that the \emph{minimum} positive weight is
$\Delta$, not that all weights are distinct.  For the mass lattice
application we need only this minimum.
\end{remark}

\section{Coordinate Rigidity}
\label{sec:rigidity}

Fix a finite generating set $\mathcal{S}$ for $\Gamma$.
The \emph{integer coordinates} of $\gamma$ are the net generator
counts $(a_1,\ldots,a_r)\in\Z^r$ in a reduced word
(Definition~\ref{def:intcoords}).

\begin{definition}\label{def:intcoords}
For generators $s_1,\ldots,s_r$, set
$a_i(\gamma)=n_i(\gamma)-n_i^-(\gamma)$, where $n_i(\gamma)$ counts
occurrences of $s_i$ and $n_i^-(\gamma)$ counts occurrences of
$s_i^{-1}$ in the geodesic word for $\gamma$.
\end{definition}

\begin{proposition}[Milnor--\v{S}varc quasi-isometry~\cite{bridsonhaefliger}]
\label{prop:msvarc}
The orbit map $\gamma\mapsto\gamma\cdot0$ is a quasi-isometry
$(\Gamma,d_\mathcal{S})\to(\Hthree,d)$.
\end{proposition}

\begin{theorem}[Coordinate rigidity]\label{thm:rigidity}
Let $\varepsilon_0=\tfrac{1}{2}\min_{f\neq f'}
d_{\mathbb{S}^2}(\xi(f),\xi(f'))$.
For any bounded deformation of magnitude $\varepsilon\leq\varepsilon_0$,
the minimal representative $\gamma_\xi$ and its integer coordinates
are unchanged.
\end{theorem}

\begin{proof}
For $\varepsilon\leq\varepsilon_0$, each perturbed direction $\xi'(f)$
lies strictly closer to $\xi(f)$ than to any other boundary direction,
so the minimal representative for $\xi'(f)$ equals that for $\xi(f)$.
\end{proof}

\section{The Arithmetic Lattice \texorpdfstring{$\pi_1(M_\mathrm{PMNS})$}{pi1(PMNS)}}
\label{sec:arithmetic}

We now specialise to $\Gamma=\pi_1(M_\mathrm{PMNS})$, where
$M_\mathrm{PMNS}$ is the closed hyperbolic 3-manifold of
second-smallest volume~\cite{Chinburg1999}, and compute the
weight-spectrum gap in terms of classical arithmetic invariants.

\subsection{The manifold and its cusped ancestor}

$M_\mathrm{PMNS}$ has volume $0.98136883$, first homology
$H_1(M_\mathrm{PMNS})=\Z/5$, and isometry group $\Z/2\times\Z/2$.
It is the $(-2,3)$ Dehn filling of the cusped manifold \texttt{m003},
which is the complement of the $(-2,3,7)$ pretzel knot in $S^3$.
The cusped manifold has volume $2.02988$ (equal to, but not isometric
with, the figure-eight knot complement~\cite{SnapPy}).
Its Chern--Simons invariant is $\mathrm{CS}(\texttt{m003})=\tfrac{1}{4}$
(exact, computed via \texttt{Sage+SnapPy}), confirming arithmeticity.

\subsection{Arithmetic invariants}

All invariants in this subsection were computed via \texttt{Sage~10+SnapPy~3.3}.

\begin{theorem}[Alexander polynomial~\cite{SnapPy}]\label{thm:alexander}
The Alexander polynomial of the $(-2,3,7)$ pretzel knot is
\[
\Delta(t) = t^2 + 3t + 1.
\]
Its roots are $-\vp^2$ and $-\vp^{-2}$, where $\vp=(1+\sqrt5)/2$.
\end{theorem}

\begin{proof}
The roots of $t^2+3t+1=0$ are $(-3\pm\sqrt5)/2$.
Since $\vp^2=(3+\sqrt5)/2$ and $\vp^{-2}=(3-\sqrt5)/2$, the roots
are $-(3+\sqrt5)/2=-\vp^2$ and $-(3-\sqrt5)/2=-\vp^{-2}$.
Verification: $\vp^2+\vp^{-2}=(\vp+\vp^{-1})^2-2=(\sqrt5)^2-2=3$.
\end{proof}

\begin{corollary}\label{cor:mahler}
The Mahler measure of $\Delta$ is $M(\Delta)=\vp^2$.
\end{corollary}

\begin{proof}
$M(\Delta)=\prod_i\max(1,|\alpha_i|)=\max(1,\vp^2)\cdot\max(1,\vp^{-2})=\vp^2$.
\end{proof}

\begin{proposition}[Adjoint torsion]\label{prop:torsion}
The adjoint torsion polynomial of \texttt{m003} is
\[
\tau_{\mathrm{adj}}(a) = -(a-1)(a^2-5a+1),
\]
with integer coefficients (verified to 29 significant figures).
\end{proposition}

\noindent This integrality is a necessary condition for $\Gamma$ to be
arithmetic~\cite{MaclachlanReid}: the adjoint torsion of an arithmetic
hyperbolic 3-manifold lies in the ring of integers of its invariant trace
field.

\begin{proposition}[Invariant trace field]\label{prop:tracefield}
The invariant trace field of \texttt{m003} is $\Q(\sqrt{-3})$.
The trace field generators satisfy $a^2+5a+7=0$ with discriminant $-3$,
confirming $\Q(\sqrt{-3})=\Q(\zeta_3)$.
\end{proposition}

\begin{remark}
The cyclotomic field $\Q(\zeta_{15})$ contains both $\Q(\sqrt{5})$
(the field of the mass lattice, regulator $\log\vp$) and
$\Q(\sqrt{-3})$ (the invariant trace field of $M_\mathrm{PMNS}$),
with $[\Q(\zeta_{15}):\Q]=\varphi(15)=8$.
The Mahler measure bridges these two subfields:
$\log M(\Delta)=2\log\vp\in\Q(\sqrt{5})\subset\Q(\zeta_{15})$.
\end{remark}

\section{The Algebraic Bridge: Lattice Spacing and Mahler Measure}
\label{sec:gap}

We now establish the algebraic connection between the mass lattice
spacing and the Mahler measure of the Alexander polynomial.

\begin{definition}[Mass lattice]
The \emph{Standard Model mass lattice} is
$\mathcal{L}=\{m_e\cdot\vp^{q/4}:q\in\Z\}$, where $m_e$ is the
electron mass.  Its logarithmic spacing is $\log\vp/4$ per unit,
giving a full-step spacing of $\log\vp=\mathrm{Reg}(\Q(\sqrt5))$.
\end{definition}

\begin{theorem}[Algebraic Bridge]\label{thm:bridge}
Let the mass lattice $\mathcal{L}=\{m_e\cdot\vp^{q/4}:q\in\Z\}$ be
given.  Its full logarithmic step spacing is
$\sigma:=\log\vp=\mathrm{Reg}(\Q(\sqrt5))$.
The Alexander polynomial of the cusped ancestor \textup{\texttt{m003}}
satisfies
\[
\tfrac{1}{2}\log M(\Delta_{\textup{\texttt{m003}}}) = \log\vp = \sigma.
\]
Hence the logarithmic spacing of the mass lattice equals half the
logarithm of the Mahler measure of the Alexander polynomial of the
cusped ancestor of $M_\mathrm{PMNS}$.
\end{theorem}

\begin{proof}
By Corollary~\ref{cor:mahler}, $M(\Delta)=\vp^2$, so
$\log M(\Delta)=2\log\vp$.  The mass lattice spacing is
$\sigma=\log\vp$ by definition.  Hence
$\sigma=\tfrac12\log M(\Delta)$.
The identification $\log\vp=\mathrm{Reg}(\Q(\sqrt5))$ holds because
$\vp$ is the fundamental unit of $\Q(\sqrt5)$ and the regulator equals
$\log|\vp|=\log\vp$.
\end{proof}

\begin{remark}
This theorem does not claim that the systole of $M_\mathrm{PMNS}$
equals $\log\vp$; the systole and the lattice spacing are different
quantities.  What the theorem asserts is a purely algebraic identity:
the lattice spacing (defined geometrically from the symmetric space
$\mathrm{SL}(2,\R)/\mathrm{SO}(2)$) coincides with half the log-Mahler
measure of the Alexander polynomial of the topological ancestor of the
manifold used to encode the lepton sector.  This coincidence is exact
and independent of any numerical fit or free parameter.
\end{remark}

\begin{corollary}[Topological origin of the golden ratio]
The golden ratio $\vp$ appears in the Standard Model mass lattice because
$M_\mathrm{PMNS}$ is the arithmetic hyperbolic 3-manifold whose
Alexander polynomial $\Delta(t)=t^2+3t+1$ has roots $-\vp^2$ and
$-\vp^{-2}$.  This identification is a theorem, not a numerical fit.
No other manifold in the volume-ordered census of closed hyperbolic
3-manifolds has an Alexander polynomial with golden-ratio roots.
\end{corollary}

\begin{remark}[Figure-eight knot]
The figure-eight knot complement (\texttt{m004}) has the same
hyperbolic volume $2.02988$ as \texttt{m003} but Alexander polynomial
$\Delta_{4_1}(t)=t-3+t^{-1}$ with Mahler measure $3+2\sqrt{2}$.
This does not give the golden ratio.  The two manifolds are not
isometric~\cite{SnapPy}; the volume coincidence is not an obstacle.
\end{remark}

\section{Discussion}

The three structural results of
Sections~\ref{sec:uniqueness}--\ref{sec:rigidity} are general
consequences of cocompactness and torsion-freeness, holding for any
arithmetic or non-arithmetic lattice.
The content specific to $M_\mathrm{PMNS}$ is
the identification of the gap constant with the Mahler measure of the
Alexander polynomial (Theorem~\ref{thm:bridge}), which exploits the
arithmeticity of $\pi_1(M_\mathrm{PMNS})$ and the specific form
$\Delta(t)=t^2+3t+1$.

The Weeks manifold $M_\mathrm{Weeks}$ (volume $0.9427$,
$H_1=\Z/5\oplus\Z/5$, isometry group $D_6$) is the $(2,1)$ Dehn
filling of the same cusped ancestor \texttt{m003}.  The Alexander
polynomial $\Delta(t)=t^2+3t+1$ and its Mahler measure $\vp^2$ are
invariants of the \emph{cusped} manifold \texttt{m003}, not of the
closed fillings.  Closed 3-manifolds do not have Alexander polynomials.
What the Weeks manifold inherits from \texttt{m003} is the arithmetic
structure: $\mathrm{CS}(\texttt{m003})=1/4$, and both the $(2,1)$ and
$(-2,3)$ fillings are arithmetic closed manifolds.
The volume hierarchy
\[
M_\mathrm{Weeks}\;(0.9427)\;<\;M_\mathrm{PMNS}\;(0.9814)\;<\;M_\mathrm{CKM}\;(2.0289)
\]
places all three flavor manifolds at the bottom of the volume spectrum,
with two of them (Weeks and PMNS) arising as Dehn fillings of the
single cusped ancestor whose Alexander polynomial has Mahler measure
$\vp^2$.  Whether the Weeks manifold plays a role in the flavor
framework beyond this arithmetic relationship is an open question.

\section{Conclusion}

We have proved three structural properties of logarithmic weight
spectra arising from cocompact torsion-free lattice actions on
$\Hthree$: uniqueness of minimal representatives, a uniform
suppression gap, and stability under bounded deformations.
For the arithmetic lattice $\pi_1(M_\mathrm{PMNS})$ we have identified
the logarithmic spacing of the associated mass lattice with
$\tfrac12\log M(\Delta_{\texttt{m003}})=\log\vp=\mathrm{Reg}(\Q(\sqrt5))$.
This provides a purely topological origin for the golden ratio in the
weight spectrum of the Standard Model mass lattice, independent of
any phenomenological fitting.

\bibliographystyle{plain}
\bibliography{gentry-hyperbolic-lattice}

\end{document}
